% ===== SECTION 9: INFORMATION GATHERING TEMPLATES (OPTIONAL) =====
\section{\texorpdfstring{\color{MidnightBlue}Information Gathering Templates (Optional)}{Information Gathering Templates (Optional)}}

\textit{``RIDE is interested in learning more about potential bidders' school bus operations and services offered. Please complete the following using the templates provided.''}

\begin{sidebar}{RFI Reference}
Section C: Information Gathering Templates (Optional) --- Technological Capabilities\\[0.3cm]
\textit{Note: Sentry's response focuses on the Technological Capabilities template, consistent with our positioning as a technology and system management provider. Staffing model and safety/compliance practices are addressed in the Company Practices Disclosure section above.}
\end{sidebar}

\subsection{Technological Capabilities}

This section describes Sentry's technology platform and its applicability to RIDE's Statewide Transportation System.

\subsubsection{Platform Overview}

Sentry's technology platform is a cloud-native, real-time transportation management system built on AWS GovCloud infrastructure. The platform is architected for multi-tenant operation---meaning a single deployment can serve multiple transportation operators, zones, and oversight authorities simultaneously, each with role-based access to their own data while RIDE maintains a unified statewide view. The system achieves 99.999\% uptime through redundant availability zones, automated failover, and continuous monitoring, ensuring that dispatch, tracking, and parent notifications are available around the clock during every school day.

The platform comprises five core modules: (1) \textbf{Routing Optimization}---AI-powered route planning with dynamic re-optimization as students enroll, withdraw, or change schools; (2) \textbf{Real-Time Dispatch \& GPS Tracking}---live vehicle monitoring with 10-second GPS update intervals, geofencing, and automated alert generation; (3) \textbf{Fleet \& Maintenance Management}---preventive maintenance scheduling, vehicle lifecycle tracking, and DOT compliance monitoring; (4) \textbf{Analytics \& Reporting}---executive dashboards, KPI scorecards, and ad hoc analysis tools; and (5) \textbf{Communication \& Notifications}---parent/guardian mobile app, SMS/email alerts, and driver in-app messaging.

Sentry's integration philosophy is vendor-neutral and API-first. The platform exposes open REST APIs and supports webhook-based event streaming, enabling integration with any GPS telematics provider, camera system, student information system, or billing platform. This is a deliberate architectural choice: RIDE should never be locked into a single hardware vendor or bus operator's proprietary system. Sentry has proven this approach at scale, integrating with 12+ third-party systems in our MTA Access-A-Ride operation, and the platform is GTFS-RT compatible for real-time transit data exchange.

\vspace{0.3cm}

RIDE has already taken steps toward technology modernization---launching a \textbf{Parent GPS App} in October~2024 and requiring operators to use digital systems for proactive outreach to LEAs and families for earlier sign-up. Sentry's platform is designed to integrate with and extend these existing investments, providing the unified backend that connects parent-facing tools, operator dispatch systems, RIDE's eRIDE student information portal, and the billing infrastructure that currently processes 700 manual checks per year from districts.

% --- Platform Architecture Diagram ---
\begin{center}
\begin{tikzpicture}[
    node distance=1.2cm and 1.5cm,
    block/.style={rectangle, draw=MidnightBlue, fill=MidnightBlue!8, text width=2.8cm, minimum height=1cm, align=center, rounded corners=3pt, font=\scriptsize},
    smallblock/.style={rectangle, draw=MidnightBlue!60, fill=MidnightBlue!4, text width=2.2cm, minimum height=0.8cm, align=center, rounded corners=2pt, font=\tiny},
    arrow/.style={-{Stealth[length=2.5mm]}, thick, MidnightBlue},
    label/.style={font=\tiny\itshape, MidnightBlue}
]
    % Data sources (left)
    \node[block] (sis) {Student\\Information\\System};
    \node[block, below=0.8cm of sis] (gps) {Vehicle GPS\\Telematics};
    \node[block, below=0.8cm of gps] (camera) {Camera \&\\Video Systems};

    % Central platform
    \node[block, right=2cm of gps, minimum height=2cm, text width=3.2cm, fill=MidnightBlue!15] (platform) {\textbf{Sentry}\\
    \textbf{Platform}\\[0.2cm]
    \tiny Routing $\cdot$ Dispatch\\
    \tiny Analytics $\cdot$ Fleet};

    % Outputs (right)
    \node[block, right=2cm of platform] (ride) {RIDE\\Dashboard};
    \node[block, above=0.8cm of ride] (parent) {Parent\\Notifications};
    \node[block, below=0.8cm of ride] (driver) {Driver\\App};

    % Below
    \node[block, below=0.8cm of driver] (violation) {Violation\\Processing};

    % Arrows
    \draw[arrow] (sis) -- node[above, label]{API} (platform);
    \draw[arrow] (gps) -- node[above, label]{Real-Time} (platform);
    \draw[arrow] (camera) -- node[above, label]{Video Feed} (platform);
    \draw[arrow] (platform) -- (ride);
    \draw[arrow] (platform) -- (parent);
    \draw[arrow] (platform) -- (driver);
    \draw[arrow] (platform) -- (violation);
\end{tikzpicture}
\end{center}

\subsubsection{Routing Optimization}

Sentry's routing engine uses AI-powered optimization algorithms to design, evaluate, and continuously refine routes across the statewide system. The engine processes student addresses, school locations, bell times, vehicle capacities, wheelchair positions, IEP-mandated maximum ride times, and real-world road network data to produce routes that minimize cost while maximizing service quality. Key capabilities include:

\begin{description}[style=nextline, leftmargin=0.5cm]
    \item[\faIcon{route} \textbf{Multi-Objective Optimization:}] The routing algorithm simultaneously minimizes total route time, deadhead miles, and ride time per student while respecting hard constraints: vehicle seating capacity, wheelchair securement positions, school bell times, maximum ride time limits specified in individual IEPs, aide assignment requirements, and RI's 18+ mile out-of-district radius. The engine evaluates millions of possible route combinations and selects the Pareto-optimal solution that balances cost efficiency against service quality. For RIDE's system, this directly targets the extreme deadhead times identified in the October 2024 system overview---particularly in Zones C (2h12m overhead) and D (3h38m overhead)---by reorganizing stop sequences and vehicle assignments to minimize empty-bus travel.

    \item[\faIcon{map-marked-alt} \textbf{Geographic Clustering:}] The engine automatically identifies clusters of students who can be efficiently served by a single route, using actual road network topology---turn restrictions, one-way streets, bridge weight limits, school zone speed reductions---rather than straight-line distance. This is critical in Rhode Island, where the state's dense road network, island geography (Block Island, Aquidneck Island), and bridge dependencies mean that proximity on a map does not equal proximity on a route. The clustering algorithm also identifies opportunities to convert bus routes to 10-passenger van routes where student density supports it, aligned with RIDE's stated intent to maximize van usage under the new legislation.

    \item[\faIcon{retweet} \textbf{Scenario Modeling:}] The platform provides robust ``what-if'' scenario modeling, enabling RIDE to evaluate the cost and service impact of policy changes before committing to them. RIDE can model alternative zone configurations---whether maintaining the current 7-zone structure, implementing the Commission's proposed 9-zone model, or testing other boundaries---and see projected route counts, deadhead reduction, vehicle requirements, and cost per student for each scenario. The engine also models van conversion opportunities (identifying which of the current 334 routes could be served by 10-passenger vans), bell-time staggering effects on route efficiency, and the impact of removing the 550 ``ghost riders'' identified by the Commission from transportation rolls.

    \item[\faIcon{sync} \textbf{Continuous Optimization:}] Routes are not static. The system continuously re-optimizes as students enroll or withdraw mid-year, schools open or close, traffic patterns shift seasonally, and new stop requests are submitted. When RIDE's eRIDE portal processes a new student enrollment or address change, the routing engine automatically evaluates whether the student can be added to an existing route or whether a route modification is needed---and presents the optimal solution to the Statewide System Manager for approval. This continuous optimization is essential for a system where ridership has grown 34\% since 2015 and close to 100 new daily routes have been added over the past decade.
\end{description}

\subsubsection{Real-Time Dispatch \& GPS Tracking}

Sentry's dispatch and tracking module provides real-time operational visibility across all zones, operators, and vehicle types from a single unified interface. In our MTA Access-A-Ride operation, this system tracks over 15,000 vehicles simultaneously---RIDE's 334-vehicle fleet is well within its proven capacity. Key capabilities include:

\begin{description}[style=nextline, leftmargin=0.5cm]
    \item[\faIcon{satellite} \textbf{Live Vehicle Tracking:}] All vehicles are tracked in real time via Geotab GO9+ GPS devices with 10-second update intervals. Dispatchers see every vehicle across all seven zones on a single interactive map, with color-coded status indicators showing on-time, delayed, off-route, and completed runs. Geofences are automatically generated around every school, stop location, and depot---triggering arrival/departure timestamps that feed into on-time performance calculations and ridership verification. The system provides RIDE with the same view that zone operators see, ensuring full transparency without relying on operator self-reporting.

    \item[\faIcon{bell} \textbf{Automated Alerts:}] The system generates automated alerts for late depot departures, route deviations beyond configurable distance thresholds, speed violations (including school zone speeding), unauthorized stops, extended idling exceeding 5 minutes, missed pickups, and vehicles approaching maximum ride-time limits for IEP-governed students. Alerts are delivered simultaneously to zone dispatchers, the Statewide System Manager, and RIDE's oversight dashboard, ensuring that service issues are identified and addressed in minutes rather than discovered after the fact through parent complaints. In our MTA operation, this alert framework achieves 95\% on-time performance across 1,000,000+ annual rides.

    \item[\faIcon{mobile-alt} \textbf{Parent Communication:}] Parents and guardians receive real-time estimated arrival notifications, geofence-triggered arrival alerts, and proactive delay notifications via mobile app, SMS, or email---based on their communication preference. RIDE has already deployed a Parent GPS App (launched October 2024); Sentry's platform is designed to integrate with and extend this existing investment via GTFS-RT data feeds, providing the real-time vehicle location backend that powers parent-facing notifications. The system also supports multilingual communication, consistent with the diverse language needs of Rhode Island's student population, and includes accessibility features compliant with WCAG 2.1 AA standards.
\end{description}

\subsubsection{Live Digital Video Violation Detection}

\mybox{
\textbf{\faIcon{video} Alignment with 2025 Legislation:} Rhode Island's Live Digital Video Violation Detection Systems act requires all new school buses to be equipped with these systems by July 1, 2027, and all buses by July 1, 2032. Sentry's video platform is designed to meet this mandate.
}{MidnightBlue}

\begin{description}[style=nextline, leftmargin=0.5cm]
    \item[\faIcon{camera} \textbf{Camera Hardware:}] Sentry's video platform supports integration with leading camera hardware manufacturers to deploy multi-angle exterior coverage on each school bus: forward-facing stop-arm cameras (capturing oncoming traffic), rear-facing cameras (capturing vehicles approaching from behind), and side-mounted cameras providing full-lane coverage. Systems capture at minimum 1080p resolution at 30 frames per second, with infrared night-vision capability and weather-sealed housings rated for New England conditions including rain, snow, ice, and road salt exposure. All camera units are tamper-proof with sealed enclosures, anti-vibration mounts, and automated health-check reporting---the system alerts maintenance staff immediately if a camera goes offline or is obstructed.

    \item[\faIcon{brain} \textbf{AI-Powered Detection:}] The system uses computer vision and machine learning algorithms to automatically detect stop-arm violations in real time. When a school bus extends its stop arm and activates flashing red lights, the AI monitors all camera feeds for vehicles that fail to stop. Upon detecting a violation, the system captures the offending vehicle's license plate using automated license plate recognition (ALPR), records a timestamped video clip showing the full violation sequence (approach, pass, and departure), and packages the evidence with GPS coordinates, date/time, bus identification, route number, and stop location. This evidence package is generated automatically within seconds of the violation---no manual video review is required at the detection stage.

    \item[\faIcon{gavel} \textbf{Violation Processing Pipeline:}] The end-to-end violation processing pipeline operates as follows: (1)~AI detection flags a potential violation and generates the evidence package; (2)~a trained human reviewer confirms the violation, verifying that the stop arm was extended, lights were flashing, and the vehicle clearly passed the bus---rejecting false positives; (3)~the system matches the license plate to DMV vehicle registration records and generates a citation; (4)~citations are filed with the appropriate jurisdiction---municipal court or the RI Traffic Tribunal---with the complete evidence package attached; (5)~fine collection is processed through the court system, with revenue distributed per the statutory formula: 30\% to the municipality, 30\% to the technology vendor, and 40\% to the state general fund (or 75\%/12.5\%/12.5\% vendor/city/state under existing Providence and Cumberland agreements). The entire workflow from violation to citation mailing targets a 14-day turnaround.

    \item[\faIcon{dollar-sign} \textbf{Revenue Potential:}] The Providence pilot program (October 2024 -- June 2025) demonstrated both the safety need and revenue potential: over 7,000 stop-arm violations were detected in just 9 months, generating approximately \$300,000 in fine revenue during FY2025---from a limited deployment on a subset of Providence routes. Scaling this capability statewide across all 334 routes, covering urban, suburban, and rural roads across all 7 zones, could reasonably generate \$1.5--2.5 million in annual violation revenue while significantly improving student safety. At \$250--500 per violation, even a modest deterrence effect---where violation rates decline as drivers become aware of the cameras---still produces substantial revenue while achieving the primary goal of protecting students at bus stops. The 30\% vendor share creates a self-funding technology model that reduces the net cost of the system management platform to RIDE and participating municipalities.
\end{description}

\subsubsection{Analytics \& Reporting}

\begin{description}[style=nextline, leftmargin=0.5cm]
    \item[\faIcon{chart-bar} \textbf{RIDE Executive Dashboard:}] The RIDE Executive Dashboard provides a real-time, single-pane-of-glass view of the entire Statewide Transportation System. RIDE leadership sees live KPIs across all zones, operators, and route types---including on-time performance by zone and operator, active vehicle count versus planned, route completion rates, average ride times versus IEP maximums, driver credential compliance rates, fleet maintenance status, and parent complaint volumes. The dashboard supports drill-down from statewide summary to individual zone, route, vehicle, or student level, enabling RIDE to identify systemic issues (e.g., a zone consistently underperforming) and isolated incidents (e.g., a specific route running late) from the same interface.

    \item[\faIcon{file-alt} \textbf{Automated Reporting:}] The platform generates automated scheduled reports at configurable intervals: daily service summaries (route completions, delays, incidents, parent communications), weekly KPI scorecards comparing each zone operator against contractual benchmarks, monthly cost analyses breaking down expenditures by zone, route type (special education, PCCT, McKinney-Vento), and cost category (driver hours, fuel, maintenance, deadhead), and annual performance reviews suitable for legislative reporting and contract renewal evaluation. All reports are delivered automatically via email to designated RIDE staff and are archived in the platform's document repository for historical trend analysis.

    \item[\faIcon{search} \textbf{Ad Hoc Analysis:}] RIDE staff have access to self-service analytics tools that allow authorized users to query the platform's data warehouse, build custom reports using drag-and-drop report builders, apply filters by date range, zone, operator, route type, or student population, and export results to Excel, CSV, and PDF formats. The system also supports saved queries and scheduled custom reports, so that frequently needed analyses (e.g., deadhead time trends by zone, ghost rider identification, cost-per-student by LEA) can be generated on demand without requiring technical support. This self-service capability ensures that RIDE is never dependent on the technology vendor or bus operators to access its own data.

    \item[\faIcon{user-shield} \textbf{FERPA Compliance:}] The platform is designed from the ground up for FERPA compliance and student data protection. All student personally identifiable information (PII) is encrypted at rest using AES-256 encryption and in transit using TLS 1.3. The system enforces role-based access controls (RBAC)---zone operators see only their assigned students, drivers see only their assigned route manifests, and RIDE administrators see the complete statewide dataset. Every data access event is logged in an immutable audit trail, recording who accessed what data, when, and from which IP address. The platform is hosted on AWS GovCloud, which meets FedRAMP High authorization standards, and complies with FERPA, HIPAA (for students with medical transportation needs), and Rhode Island's student privacy statutes. Data retention policies are configurable to meet RIDE's records management requirements, and all student data can be purged upon contract termination.
\end{description}

\subsubsection{Integration Capabilities}

\begin{tabularx}{\textwidth}{|X|X|X|}
\hline
\rowcolor{tableHeader}
\textcolor{white}{\textbf{System}} &
\textcolor{white}{\textbf{Integration Method}} &
\textcolor{white}{\textbf{Data Exchanged}} \\
\hline
eRIDE (RIDE's Student Portal) & API / Flat File & Enrollment, addresses, IEP transportation plans, school assignments, annual sign-up data \\ \hline
RIDE Parent GPS App & API / GTFS-RT & Real-time vehicle location, ETA, arrival notifications (extending RIDE's existing October~2024 deployment) \\ \hline
GPS / Telematics Providers & API / GTFS-RT & Vehicle location, speed, engine diagnostics \\ \hline
Camera / Video Systems & API / SDK & Video feeds, violation events, evidence packages \\ \hline
LEA Billing / State Aid & API / Export & Automated ``expenditure follows the child'' cost allocation, replacing 700 annual manual checks with real-time per-student billing \\ \hline
RIDE Reporting Systems & API / Dashboard & KPIs, compliance reports, audit data \\ \hline
\end{tabularx}
